\section{Constraints}

This section fixes the SHA application matrices.  SHA-256 word semantics are
those of FIPS~180-4.  The row templates and the one-compression counts are
those proved in
\href{https://github.com/reilabs/freigen/commit/93d3c55d69b4e7747eacf24212f93c795d05fdba}
{Freigen commit \code{93d3c55}}; this profile additionally fixes public-word
binding and the canonical row sharing for a chain of $N$ compressions.

\subsection{Statement and compact matrix geometry}

The public statement contains the initial and final chaining values and all
message blocks:
\begin{equation}
 x_{\rm sha}:=(N,H_0,(B_i)_{i=0}^{N-1},H_N),\qquad
 H_i\in(\{0,1\}^{32})^8,\quad B_i\in(\{0,1\}^{32})^{16}.
 \label{eq:sha-statement}
\end{equation}
Every word in \cref{eq:sha-statement} has big-endian byte encoding.  The
canonical $\mathsf{app\_statement\_bytes}$ are $\operatorname{LE}_{32}(N)$
followed by $H_0,B_0,\ldots,B_{N-1},H_N$, four big-endian bytes per word and
no inner lengths.  Their exact length is $68+64N$ bytes; trailing bytes reject.
This profile requires
$1\le N\le429389$, where the upper endpoint is the largest $N$ for which
$20005N+257\le2^{33}$.  Before allocating, the statement preflight computes
$68+64N$, $n_f$, $n_h$, and $n_A$ with checked integer arithmetic and rejects
on any overflow or bound failure.  Padding and private-message hashing require
different application profiles.

The compact Boolean input and virtual-output lengths, excluding the affine
constant coordinate, are
\begin{equation}
 n_f:=6641N+256,\qquad n_h:=20005N+256.
 \label{eq:sha-live-lengths}
\end{equation}
Set
\begin{equation}
 \begin{aligned}
 m_f&:=\max(22,\lceil\log_2n_f\rceil),&
 m_h&:=\max(22,\lceil\log_2(n_h+1)\rceil),\\
 (t_f,s,W_{\rm PCS})&=(m_f-7,7,1),&
 (t_h,s,W_{\rm PCS})&=(m_h-7,7,1).
 \end{aligned}
 \label{eq:sha-shapes}
\end{equation}
The profile requires $22\le m_f\le m_h\le33$.  Otherwise it rejects and a
different profile is required.  In particular the virtual chunk width is
$127-t_h-W_{\rm PCS}=133-m_h\ge100$, so $L=1$ and
$\mathcal A=\{0\}$.  The packed-opening domain exponent is $m_p=t_f$.

Let $f\in\Ftwo^{n_f}$ be ordered as follows: all public message bits, the
initial-state bits, and then the hinted decomposition bits in circuit execution
order.  Within a word, bit $j=0$ is least significant and
\begin{equation}
 [X]:=\sum_{j=0}^{31}2^jX[j],\qquad
 \operatorname{lo}_{32}(Q)[j]:=Q[j]\quad(0\le j<32).
 \label{eq:sha-word-value}
\end{equation}
The execution order within each compression is: rounds $0,\ldots,15$;
schedule-plus-round steps $16,\ldots,61$; inlined round $62$; and the fused
round-$63$ feed-forward.  Adjacent compressions reuse the previous output bits
as the next state rather than allocating a duplicate state.

Let $D_f\in\Ftwo^{2^{t_f}\times2^s}$ be the committed tensor and let $J_f$
inject the first $n_f$ row-major coordinates into that tensor, so its transpose
selects them and
$f=J_f^{\mathsf T}\operatorname{vec}(D_f)$.  The remaining committed padding
is ignored by this relation; an honest encoder zeros it.  The inherited PCS
packing order is the transpose of this semantic order.  Define
\begin{equation}
 \begin{aligned}
 B_f[c,b]&:=D_f[b,c],\\
 [\Pi_f\operatorname{vec}(D_f)]_{c2^{t_f}+b}
   &:=[\operatorname{vec}(D_f)]_{b2^s+c},
   \qquad\operatorname{vec}(B_f)=\Pi_f\operatorname{vec}(D_f),
 \end{aligned}
 \label{eq:sha-pack-permutation}
\end{equation}
for $0\le b<2^{t_f}$ and $0\le c<2^s$.  Thus the committed object is
\begin{equation}
 O:=\Enc(\Pack(B_f)),\qquad \rho:=\Root(O).
 \label{eq:sha-commitment}
\end{equation}
Equation~\eqref{eq:sha-pack-permutation} is the only semantic-to-PCS index
permutation.  P8/V8 applies it to the adjoint coefficient, and P10/V10 uses
the resulting PCS order for the ring switch.

\subsection{The binary matrix $M_{\rm SHA}$}

Use augmented vectors $\bar f=(1,f)$ and $\bar h=(1,h_{\rm sem})$.  The
Freigen matrix is the fixed binary matrix
\begin{equation}
 M_{\rm SHA}\in\Ftwo^{(n_h+1)\times(n_f+1)},\qquad
 \bar h=M_{\rm SHA}\bar f.
 \label{eq:sha-m-matrix}
\end{equation}
Let $J_h$ inject $\bar h$ into the first $n_h+1$ coordinates of
$D_h\in\Ftwo^{2^{t_h}\times2^s}$.  Split
$M_{\rm SHA}=[m_0\mid M_1]$ into its constant and nonconstant columns.  The
PCS form is
\begin{equation}
 \operatorname{vec}(D_h)=h_0+T\operatorname{vec}(D_f),\quad
 h_0:=J_hm_0,\quad T:=J_hM_1J_f^{\mathsf T}.
 \label{eq:sha-virtual-map}
\end{equation}
Only $D_f$ is committed.  Neither $M_{\rm SHA}$ nor $T$ is materialized:
their rows are emitted from the following parity templates.

For a word $X$ and $0\le j<32$, define
\begin{equation}
 \begin{aligned}
 \Sigma_0(X)[j]&=X[j+2]\oplus X[j+13]\oplus X[j+22],\\
 \Sigma_1(X)[j]&=X[j+6]\oplus X[j+11]\oplus X[j+25],\\
 \sigma_0(X)[j]&=X[j+7]\oplus X[j+18]\oplus X[j+3],\\
 \sigma_1(X)[j]&=X[j+17]\oplus X[j+19]\oplus X[j+10],
 \end{aligned}
 \label{eq:sha-m-rotations}
\end{equation}
where every rotation index is modulo $32$, while the final
$\sigma_0,\sigma_1$ shift term is zero when
$j+3\ge32$ or $j+10\ge32$, respectively.  The other nonlinear-function
characters are
\begin{equation}
 P_{ef}=e\oplus f,\qquad P_{eg}=e\oplus g,\qquad
 P_{abc}=a\oplus b\oplus c.
 \label{eq:sha-m-parities}
\end{equation}
The integer combinations used below are consequently
\begin{equation}
 C_h:=[f]+[g]-[P_{ef}]+[P_{eg}]=2\Ch(e,f,g),\qquad
 C_m:=[a]+[b]+[c]-[P_{abc}]=2\Maj(a,b,c).
 \label{eq:sha-chmaj}
\end{equation}

The remaining rows of $M_{\rm SHA}$ are identity rows for input and hinted
bits.  A 34-bit schedule result derives bit zero as the XOR of its four
addends' low bits and commits only bits $1,\ldots,33$; a 33-bit ordinary
feed-forward result does the same for its two addends.  All bits of the
$Q^e$ and $Q^a$ decompositions are committed.  Concretely, for
$16\le k\le61$ and $j\in\{1,2,3,5,6,7\}$,
\begin{equation}
 \begin{aligned}
 Q^W_k[0]&=W_{k-16}[0]\oplus\sigma_0(W_{k-15})[0]
              \oplus W_{k-7}[0]\oplus\sigma_1(W_{k-2})[0],\\
 Q^{\rm out}_j[0]&=H_i[j][0]\oplus S_{64}[j][0].
 \end{aligned}
 \label{eq:sha-derived-low-bits}
\end{equation}
The two identities in \cref{eq:sha-derived-low-bits} are $M_{\rm SHA}$ rows,
not R1CS rows.  In canonical row order the exact census is
\begin{equation}
\begin{array}{@{}lr@{}}
\toprule
\text{row family in }M_{\rm SHA} & \text{rows}\\
\midrule
\text{affine constant} & 1\\
\text{message inputs; initial state} & 512N+256\\
\Sigma_0,\Sigma_1,P_{ef},P_{eg},P_{abc}\text{ in 64 rounds} & 10240N\\
\sigma_0,\sigma_1\text{ in 48 schedule steps} & 3072N\\
Q^W_{16},\ldots,Q^W_{61} & 1564N\\
Q^e\text{ decompositions} & 2242N\\
Q^a\text{ decompositions} & 2177N\\
\text{six remaining feed-forward decompositions} & 198N\\
\midrule
\textbf{total} & \mathbf{20005N+257}
\\\bottomrule
\end{array}
 \label{eq:sha-m-row-census}
\end{equation}
The augmented column count, in the same constant-first convention, is
\begin{equation}
 1+(512N+256)+1518N+2242N+2177N+192N
   =6641N+257.
 \label{eq:sha-m-column-census}
\end{equation}
For $N=1$, \cref{eq:sha-m-row-census,eq:sha-m-column-census} give exactly
$20262$ rows and $6898$ columns.  The previous $20457\times7145$ system loses
$64$ rows/columns by coupling $a^+$ to $e^+$, $66$ by inlining $W_{62},W_{63}$,
$65$ by fusing final feed-forward, and another $52$ columns by deriving low
sum bits.

\subsection{The linear $\mathbb Z$-R1CS rows}

For compression $i$, let $(a_k,b_k,c_k,d_k,e_k,f_k,g_k,h_k)$ denote the
state before round $k$, and let $W_k$ denote its schedule word.  For
$0\le k\le61$, the two round rows are
\begin{equation}
 \begin{aligned}
 2[Q^e_k]&=2\bigl([d_k]+[h_k]+[\Sigma_1(e_k)]+K_k+[W_k]\bigr)+C_h(k),\\
 2[Q^a_k]&=2\bigl([e_{k+1}]+[\Sigma_0(a_k)]-[d_k]\bigr)+C_m(k)+2^{33},
 \end{aligned}
 \label{eq:sha-r1cs-round}
\end{equation}
where $Q^e_k$ is 35 bits, $Q^a_k$ is 34 bits,
$e_{k+1}=\operatorname{lo}_{32}(Q^e_k)$, and
$a_{k+1}=\operatorname{lo}_{32}(Q^a_k)$.  The other six words are wiring
aliases, not R1CS rows:
$(b_{k+1},c_{k+1},d_{k+1},f_{k+1},g_{k+1},h_{k+1})
=(a_k,b_k,c_k,e_k,f_k,g_k)$.
Equation~\eqref{eq:sha-r1cs-round} uses
$a^+=e^+-d+\Sigma_0(a)+\Maj(a,b,c)\pmod{2^{32}}$; it does not allocate a
separate $T_1$ or 35th $Q^a$ bit.

For $16\le k\le61$, one schedule row is also present:
\begin{equation}
 [Q^W_k]=[W_{k-16}]+[\sigma_0(W_{k-15})]+[W_{k-7}]+[\sigma_1(W_{k-2})],
 \qquad W_k=\operatorname{lo}_{32}(Q^W_k),
 \label{eq:sha-r1cs-schedule}
\end{equation}
where $Q^W_k$ is 34 bits and its low bit is derived as specified above.
Define the unmaterialized terminal sums
\begin{equation}
 S_k:=[W_{k-16}]+[\sigma_0(W_{k-15})]+[W_{k-7}]+[\sigma_1(W_{k-2})]
 \quad(k\in\{62,63\}).
 \label{eq:sha-terminal-schedule}
\end{equation}
Round $62$ has only
\begin{equation}
 \begin{aligned}
 2[Q^e_{62}]&=2\bigl([d_{62}]+[h_{62}]+[\Sigma_1(e_{62})]+K_{62}+S_{62}\bigr)
   +C_h(62),\\
 2[Q^a_{62}]&=2\bigl([e_{63}]+[\Sigma_0(a_{62})]-[d_{62}]\bigr)
   +C_m(62)+2^{33},
 \end{aligned}
 \label{eq:sha-r1cs-round62}
\end{equation}
with 36-bit $Q^e_{62}$ and 34-bit $Q^a_{62}$.

Round $63$ is fused with the first and fifth feed-forward words.  If
$A_{\rm in}=H_i[0]$ and $E_{\rm in}=H_i[4]$, its rows are
\begin{equation}
 \begin{aligned}
 2[Q^{\rm out}_e]&=2\bigl([d_{63}]+[h_{63}]+[\Sigma_1(e_{63})]+K_{63}
       +S_{63}+E_{\rm in}\bigr)+C_h(63),\\
 2[Q^{\rm out}_a]&=2\bigl([\operatorname{lo}_{32}(Q^{\rm out}_e)]+[\Sigma_0(a_{63})]
       +A_{\rm in}-E_{\rm in}-[d_{63}]\bigr)+C_m(63)+2^{34}.
 \end{aligned}
 \label{eq:sha-r1cs-final}
\end{equation}
Here $Q^{\rm out}_e$ is 36 bits and $Q^{\rm out}_a$ is 35 bits; their low
32 bits are $H_{i+1}[4]$ and $H_{i+1}[0]$.  For
$j\in\{1,2,3,5,6,7\}$, the six remaining rows are
\begin{equation}
 [Q^{\rm out}_j]=H_i[j]+S_{64}[j],\qquad
 H_{i+1}[j]=\operatorname{lo}_{32}(Q^{\rm out}_j),
 \label{eq:sha-r1cs-feedforward}
\end{equation}
where $S_{64}$ is the final pre-feed-forward state and each
$Q^{\rm out}_j$ is 33 bits with a derived low bit.

The internal row count per compression is therefore
\begin{equation}
 62\cdot2+46+2+2+6=180.
 \label{eq:sha-r1cs-census}
\end{equation}
This application additionally binds the sixteen public message words per
compression and the eight words of each public endpoint.  Intermediate states
are shared.  Set $n_A:=196N+16$ and
$d_A:=\lceil\log_2n_A\rceil$.  With
$h:=\lift(\operatorname{vec}(D_h))\in\mathbb Z^{2^{m_h}}$, the row-padded
integer constraint matrix is
\begin{equation}
 \begin{aligned}
 A_x&\in\mathbb Z^{2^{d_A}\times2^{m_h}},& A_xh&=0,\\
 A_x[\beta,*]&=0&& (n_A\le\beta<2^{d_A}).
 \end{aligned}
 \label{eq:sha-relation-census}
\end{equation}
Rows are ordered blockwise as
\cref{eq:sha-r1cs-round,eq:sha-r1cs-schedule,%
eq:sha-r1cs-round62,eq:sha-r1cs-final,eq:sha-r1cs-feedforward}, followed by
message pins and the two endpoint-pin families.  Omitted columns of $A_x$ are
zero; an implementation stores row templates, not the padded matrix.

Finally, projection to the fixed PCS field is accepted only after the verifier
recomputes
\begin{equation}
 B_x:=\max_{\beta<2^{d_A}}\sum_z|A_x[\beta,z]|,\qquad
 2B_x<2^{39}<q=2^{100}-15.
 \label{eq:sha-residual-bound}
\end{equation}
Equation~\eqref{eq:sha-residual-bound} makes reduction modulo $q$ injective on
every reachable bit-witness residual.  Failure is a hard reject.
